\title{xLSTM: Extended Long Short-Term Memory}

\begin{abstract}
During the 1990s, foundational innovations in Long Short-Term Memory (LSTM) arose through the concepts of the constant error carousel and gating mechanisms. LSTMs subsequently became a pivotal technology within deep learning, serving as the backbone for the initial wave of Large Language Models (LLMs). However, the introduction of the Transformer architecture—with its emphasis on parallelizable self-attention—ushered in a transformative shift that eventually eclipsed LSTMs in scalability and performance. In this work, we pose a fundamental question: To what extent can LSTMs be advanced in language modeling if they are scaled to billions of parameters using up-to-date methods from current LLMs, all the while addressing recognized drawbacks inherent in traditional LSTM models? Our approach begins with the proposal of exponential gating, enhanced by rigorous normalization and stabilization methods. We also re-envision LSTM memory configurations, leading to: (i) sLSTM, which employs a scalar-valued memory and update, alongside a novel approach to memory mixing, and (ii) mLSTM, a variant with matrix memory and covariance-based updates that enables complete parallelism. By incorporating these advancements within residual block architectures, we construct xLSTM blocks, which can be composed into deeper xLSTM model stacks. The introduction of exponential gating and revised memory strategies significantly strengthens xLSTM’s modeling power, allowing it to match and sometimes surpass the contemporary Transformers and State Space Models with regard to both efficiency and scalability.\\
Code available at: \url{https://github.com/NX-AI/xlstm}
\end{abstract}

\section{Introduction}
\label{sec:intro}

The concepts behind Long Short-Term Memory (LSTM)~\citep{Hochreiter:91,Hochreiter:97nips,Hochreiter:97}, specifically the constant error carousel and various gating mechanisms, were proposed as solutions to the vanishing gradient issue that occurs in recurrent neural networks~\citep{Hochreiter:91,Hochreiter:00book}:

\begin{align}
\label{eq:lstm_idea}
  c_t \ = \ f_t \ c_{t-1} \ + \ 
          i_t \  z_t \ , \quad  
          h_t \ = \ o_t \ \psi({c_t}) \ . 
\end{align}

The constant error carousel operates by incrementally updating the cell state $c_{t-1}$ (indicated in green) using cell inputs $z_t$, with the update process regulated by sigmoid-based gates (depicted in blue). Specifically, the input gate $i_t$ and forget gate $f_t$ determine how the cell state is modified, whereas the output gate $o_t$ governs the cell’s contribution to the hidden state $h_t$. The cell state is first normalized or compressed via the function $\psi$, and then the application of the output gate produces the hidden state.

LSTMs have found widespread use across multiple application areas \citep{Hochreiter:01,Hochreiter:07,Schmidhuber:15}, dominating text generation methods up until the emergence of Transformers in 2017~\citep{Vaswani:17}. Their capabilities have been demonstrated in an array of sequential tasks, including but not limited to text generation~\citep{Graves:13,Karpathy:15blog}, handwriting synthesis~\citep{Graves:13}, sequence-to-sequence machine translation~\citep{Sutskever:14nips}, computer program evaluation~\citep{Zaremba:14arxiv}, producing image captions~\citep{Karpathy:15,Hossain:19}, source code generation~\citep{Karpathy:15blog}, rainfall-runoff forecasting~\citep{Kratzert:18,Kratzert:19}, and the development of hydrological models for flood forecasting~\citep{Nearing:24}. In reinforcement learning, LSTMs remain among the most effective sequence modeling architectures, serving as the backbone for systems such as AlphaStar for StarCraft II~\citep{Vinyals:17short}, OpenAI Five for Dota 2~\citep{OpenAI:19}, and controllers for magnetic confinement in nuclear fusion experiments~\citep{Degrave:22short}. Their strength lies in the ability to capture high-level abstractions, extracting and storing semantic features within their memory cells~\citep{Karpathy:15blog}; this property is illustrated by the presence of number and syntax neurons~\citep{Lakretz:19}, linguistic neurons~\citep{Bau:19}, and sentiment neurons~\citep{Radford:17arxiv}. LSTMs continue to see adoption in contemporary, high-stakes applications~\citep{Degrave:22short,Nearing:24}, reflecting their enduring utility and robustness.

Although LSTMs have achieved remarkable results, they present three principal drawbacks: (i) Inability to amend previously made storage choices. This issue is illustrated through the 
{\it Nearest Neighbor Search} task (additional details in Appendix~\ref{sec:appNearestNeighborSearch}): 
Given a reference vector, the system must iterate over a sequence to identify and retrieve the value paired with the most similar vector, which is determined only at the end of the sequence. The left panel of Figure~\ref{fig:lstmProblems} presents the mean squared error on this task. LSTMs encounter difficulties updating stored values upon encountering a superior match later in the sequence; however, our proposed xLSTM addresses this deficit via the use of exponential gating. (ii) Restricted storage capabilities, where all information must be condensed into single scalar cell states. This constraint is demonstrated using {\it Rare Token Prediction}. The right panel of Figure~\ref{fig:lstmProblems} depicts the perplexity scores for token prediction across varying frequency bins on Wikitext-103~\citep{Merity:17}. LSTMs are less effective at predicting infrequent tokens as a result of their limited memory capacity. In contrast, our xLSTM leverages a matrix-based memory to overcome this limitation. (iii) Poor parallelization due to memory entanglement, arising from the dependencies between hidden states at consecutive time steps, thus necessitating purely sequential computation.

The shortcomings of LSTM networks have led to the rise of Transformers~\citep{Vaswani:17} as a preferred approach in language modeling. A natural question then is: if we address these weaknesses and expand LSTMs to the scale of present-day Large Language Models, what levels of performance might be attainable in language modeling tasks?

\section{Extended Long Short-Term Memory}
\label{sec:xLSTM}

In response to the constraints inherent in LSTM, the Extended Long Short-Term Memory (xLSTM) framework presents two key advancements over the concept outlined in Equation~\eqref{eq:lstm_idea}. The enhancements—namely, exponential gating mechanisms and innovative memory architectures—add two distinct variants to the LSTM lineage: (i) sLSTM (detailed in Section~\ref{sec:sLSTM}), which utilizes a scalar memory, a scalar update operation, and introduces a form of memory mixing; and (ii) mLSTM (explored in Section~\ref{sec:mLSTM}), which is designed around a matrix-based memory alongside an outer product covariance update, facilitating complete parallel computation. Both sLSTM and mLSTM leverage exponential gating for improved functionality. In pursuit of parallel execution, mLSTM dispenses with memory mixing—specifically, it removes the hidden-to-hidden recurrent links. The architectures for both sLSTM and mLSTM permit the integration of multiple memory cells; for sLSTM, this includes memory mixing between the cells. Moreover, sLSTM accommodates multiple heads with memory mixing restricted to intra-head cell interactions, thus not allowing mixing between heads themselves. Incorporating heads into sLSTM, coupled with exponential gating, introduces a novel form of memory interaction. For mLSTM, configurations employing multiple heads are functionally equivalent to those using multiple cells.

When these alternative LSTM formulations are incorporated into residual block structures, the resulting modules are termed xLSTM blocks (refer to Section~\ref{sec:xLSTMblock}). Building models by stacking these xLSTM blocks in a residual manner yields xLSTM architectures, as described in Section~\ref{sec:xLSTMarchitecure}. For an illustration of the overall xLSTM design and its building blocks, see Figure~\ref{fig:xlstm_sketch}.

\subsection{Review of the Long Short-Term Memory}
\label{sec:orgLSTM}

The foundational concept behind LSTM networks~\citep{Hochreiter:91,Hochreiter:97nips,Hochreiter:97} incorporated a scalar memory cell serving as the core mechanism for information processing and retention. This design effectively addresses the vanishing gradient problem~\citep{Hochreiter:91,Hochreiter:00book} through the constant error carousel, which governs the evolution of the cell state. The structure of the memory cell comprises three distinct gates: the input gate, the output gate, and the forget gate—the latter of which was subsequently added by \citet{Gers:00}. At each time step $t$, the dynamics of the LSTM memory cell are governed by the following update equations:

\begin{align}
c_t \ &= \  f_t \ c_{t-1} \ + \ i_t \ z_t &  & &\text{cell state} \\
h_t  \ &= \ o_t \ \tilde{h}_t \ , & \tilde{h}_t \ &= \ \psi \left( c_t\right)
  &\text{hidden state} \\
z_t \ &= \ \varphi \left( \tilde{z}_t \right) \ , 
  &\tilde{z}_t \ &=  \ \mathbf{w}^\top_{z} \ \mathbf{x}_t \ + \
  r_{z}  h_{t-1} \ + \  b_{z} \ \
  &\text{cell input} \\
i_t \ &= \ \sigma \left( \tilde{i}_t  \right) \ , 
  &\tilde{i}_t \ &= \ \mathbf{w}^\top_{i} \ \mathbf{x}_t \ + \
  r_{i}  \ h_{t-1} \ + \  b_{i} \ \
  &\text{input gate} \\
f_t \ &= \ \sigma \left(  \tilde{f}_t \right) \ , 
  &\tilde{f}_t \ &= \ \mathbf{w}^\top_{f} \ \mathbf{x}_t  \ + \
  r_{f}  \ h_{t-1} \ + \  b_{f} \ \
  &\text{forget gate} \\
o_t \ &= \ \sigma \left( \tilde{o}_t \right) \ , 
  &\tilde{o}_t  \ &= \ \mathbf{w}^\top_{o} \ \mathbf{x}_t \ + \
  r_{o}  \ h_{t-1} \ + \  b_{o} \ \
  &\text{output gate} 
\end{align}

The input weight vectors $\mathbf{w}_{z}$, $\mathbf{w}_{i}$, $\mathbf{w}_{f}$, and $\mathbf{w}_{o}$ map the external input $\mathbf{x}_t$ to the cell input, input gate, forget gate, and output gate, respectively. In parallel, the recurrent weights $r_{z}$, $r_{i}$, $r_{f}$, and $r_{o}$ connect the hidden state $h_{t-1}$ to these same components: cell input, input gate, forget gate, and output gate. Each of these pathways includes its own bias parameter, specifically $b_{z}$, $b_{i}$, $b_{f}$, and $b_{o}$. The nonlinearities $\varphi$ and $\psi$ act as activation functions for the cell input and hidden state, often realized as $\tanh$. In particular, $\psi$ plays a role in constraining the cell state, preventing it from diverging. Gate activations utilize the sigmoid function, defined as $\sigma \left( x \right)= 1/(1+\exp(-x))$. Subsequent developments introduced vector-valued memory states $\mathbf{c}_t \in \mathbb{R}^{d}$, facilitating the application of recurrent weight matrices $\mathbf{R} \in \mathbb{R}^{d \times d}$, which enable mixing between outputs of different memory cells~\citep{Greff:15}; further explanation is provided in Appendix~\ref{sec:apporgLSTM}. Investigations by ablation have indicated the necessity of preserving all elements within the memory cell structure~\citep{Greff:15}.

\subsection{sLSTM}
\label{sec:sLSTM}

To enable LSTMs to modify their storage choices, we propose the integration of exponential gates (highlighted in red) alongside mechanisms for normalization and stabilization. Specifically, we allow the input and forget gates to utilize exponential activation functions. Regarding normalization, we define a normalizer state that accumulates the sum of the input gate multiplied by all subsequent forget gates.

The scalar sLSTM forward pass is:

\begin{align}
\label{eq:slstmforward}
c_t \ &= \  f_t \ c_{t-1} \ + \ i_t \ z_t & & 
 &\text{cell state} \\
n_t \ &= \  f_t \ n_{t-1} \ + \ i_t  & & 
 &\text{normalizer state} \\
h_t \ &= \ o_t \ \tilde{h}_t \ ,
  &\tilde{h}_t \ &= \ \frac{c_t}{n_t}
&\text{hidden state} \\
z_t \ &= \ \varphi \left( \tilde{z}_t \right) \ , 
  &\tilde{z}_t \ &=  \ \mathbf{w}^\top_{z} \ \mathbf{x}_t \ + \
  r_{z}  h_{t-1} \ + \  b_{z}
  &\text{cell input} \\
i_t \ &= \ \!\exp \left( \tilde{i}_t  \right) \ , 
  &\tilde{i}_t \ &= \ \mathbf{w}^\top_{i} \ \mathbf{x}_t \ + \
  r_{i}  \ h_{t-1} \ + \  b_{i}
  &\text{input gate} \\
f_t \ &= \ \sigma \left(  \tilde{f}_t \right) \ \text{OR} \
\!\exp \left(  \tilde{f}_t \right) \ ,
  &\tilde{f}_t \ &= \ \mathbf{w}^\top_{f} \ \mathbf{x}_t  \ + \
  r_{f}  \ h_{t-1} \ + \  b_{f}
  &\text{forget gate} \\
o_t \ &= \ \sigma \left( \tilde{o}_t \right) \ , 
  &\tilde{o}_t  \ &= \ \mathbf{w}^\top_{o} \ \mathbf{x}_t \ + \
  r_{o}  \ h_{t-1} \ + \  b_{o}
  &\text{output gate} 
\end{align}

We adapt the foundational gating mechanisms from LSTMs—including input- and hidden-dependent gates along with a bias term—into the proposed architectures. Because exponential activations may result in excessively large outputs and potential numerical overflows, we introduce an auxiliary state \mbox{$m_t$~\citep{Milakov:18arxiv}} to provide stability to the gates:

\begin{align}
\label{eq:slstmstabil}
m_t \ &= \ \max \left( \log ( f_t ) + m_{t-1} , \log ( i_t ) \right) &\text{stabilizer state} \\
i'_t \ &= \ \!\exp \left( \log \left ( i_t \right) - m_t \right) \ = \ \!\exp \left( \tilde{i}_t - m_t \right) \ \, 
  &\text{stabil. input gate} \\
f'_t \ &= \ \!\exp \left( \log \left( f_t \right) + m_{t-1} - m_t \right) \ \, 
  &\text{stabil. forget gate}
\end{align}

As demonstrated in Appendix~\ref{sec:appsLSTM}, substituting $f_t$ with $f'_t$ and $i_t$ with $i'_t$ during the forward computation leaves both the network's overall output and the loss gradients with respect to the parameters unchanged.

\paragraph{New Memory Mixing.} Similar to the classical LSTM, sLSTM supports multiple memory cells (refer to Appendix~\ref{sec:appsLSTM}). When there are multiple memory cells, they facilitate memory mixing through the recurrent matrices $\mathbf{R}_{\mathbf{z}}$, $\mathbf{R}_{\mathbf{i}}$, $\mathbf{R}_{\mathbf{f}}$, $\mathbf{R}_{\mathbf{o}}$, which connect the hidden state $\mathbf{h}$ to the cell input $\mathbf{z}$ and the gating units $\mathbf{i}$, $\mathbf{f}$, and $\mathbf{o}$. A distinguishing factor in this architecture is the incorporation of exponential gating during memory mixing. In the latest sLSTM design, each head can include several memory cells interacting via mixing; however, such interactions are confined within the same head and do not span across different heads. The use of distinct heads in sLSTM, combined with exponential gating, creates a novel approach to memory mixing.

\subsection{mLSTM}
\label{sec:mLSTM}

To improve the capacity of LSTM models, we modify the traditional scalar memory cell $c \in \mathbb{R}$ to a matrix form $\mathbf{C} \in \mathbb{R}^{d \times d}$. This allows retrieval operations to utilize matrix multiplication. At each time step $t$, the architecture stores two vectors: a key $\mathbf{k}_t \in \mathbb{R}^d$ and a value $\mathbf{v}_t \in \mathbb{R}^d$ (terminology borrowed from the Transformer model). At a future time $t + \tau$, we aim to recover the stored value $\mathbf{v}_t$ by employing a query vector $\mathbf{q}_{t+\tau} \in \mathbb{R}^d$. This methodology corresponds to the framework known as Bidirectional Associative Memories (BAMs) \citep{Kohonen:72,Anderson:72,Nakano:72,Anderson:77}.

The covariance update rule~\citep{Sejnowski:77,Dayan:91} for storing a key-value pair is

\begin{align}
\mathbf{C}_t \ &= \ \mathbf{C}_{t-1} \ + \ \mathbf{v}_{t} \ \mathbf{k}_t^\top \ .
\end{align}

We posit that inputs undergo layer normalization prior to their projection into key and value representations, ensuring a zero-mean for these components. The process of updating the covariance matrix is theoretically optimal~\citep{Dayan:91} when aiming for maximal discrimination among retrieved binary vectors, a condition that also corresponds to optimizing the signal-to-noise ratio. Permitting only pairwise retrievals—akin to the mechanism of attention with quadratic computational cost—enables further enhancement of vector separability~\citep{Krotov:16,Krotov:17,Ramsauer:21}. Notably, the covariance-based update mechanism aligns with that utilized in Fast Weight Programmers~\citep{Schmidhuber:92ncfastweights,Schlag:21}, later refined by the introduction of a fixed decay factor for $\mathbf{C}_{t-1}$ and a fixed learning rate for the term 
$\mathbf{v}_{t} \mathbf{k}_t ^\top$~\citep{Ba:16}. Consistent with these developments, we adapt the covariance update rule within the LSTM architecture, associating the forget gate with the decay rate, the input gate with the learning rate, and using the output gate to modulate the resulting retrieved vector.

Within this matrix memory framework, the normalizer state represents a weighted aggregation of key vectors; each key vector's weight is determined by both the input gate and the product of all subsequent forget gates. The normalizer state, as before, tracks the cumulative influence of these gates. Notably, since the dot product between the query and the normalizer state may approach zero, we adopt the approach of taking its absolute value and enforcing a minimum threshold (usually 1.0), consistent with previous practice~\citep{Sun:23arxiv}. The forward computation of the mLSTM proceeds as follows:

\begin{align}
\label{eq:mlstm_recurrent_begin}
\mathbf{C}_t \ &= \  f_t \ \mathbf{C}_{t-1} \ + \ 
  i_t \ \mathbf{v}_t \ \mathbf{k}_t^\top & &  &\text{cell state} \\
\mathbf{n}_t \ &= \  f_t \ \mathbf{n}_{t-1} \ + \ 
  i_t \ \mathbf{k}_t & &  &\text{normalizer state} \\
\mathbf{h}_t  \ &= \ \mathbf{o}_t \ \odot \ \tilde{\mathbf{h}}_t
\ , \qquad \qquad 
\tilde{\mathbf{h}}_t \ = \   \mathbf{C}_t \mathbf{q}_t \ / \ 
  \max \left\{ |\mathbf{n}_t^\top \mathbf{q}_t|, 1 \right\} 
   &&&\text{hidden state} \\
\mathbf{q}_t \ &= \ \mathbf{W}_q \ \mathbf{x}_t \ + \ \mathbf{b}_q  & &  &\text{query input} \\
\mathbf{k}_t \ &= \ \frac{1}{\sqrt{d}} \mathbf{W}_k \ \mathbf{x}_t \ + \ \mathbf{b}_k  & &  &\text{key input} \\
\mathbf{v}_t \ &= \ \mathbf{W}_v \ \mathbf{x}_t \ + \ \mathbf{b}_v  & &  &\text{value input} \\
i_t \ &= \ \!\exp \!  \! \left( \tilde{i}_t  \right) \ , 
 \qquad \qquad \ \ \ \ \,
  \tilde{i}_t \ = \ \mathbf{w}^\top_{i} \ \mathbf{x}_t \ + \  b_{i} \ \ \ 
  &&&\text{input gate} \\
f_t \ &= \ \sigma \!  \left(  \tilde{f}_t \right) \ \text{OR} \ \!\exp \!  \! \left(  \tilde{f}_t \right) , \ \ \: \!
\tilde{f}_t \ = \ \mathbf{w}^\top_{f} \ \mathbf{x}_t  \ + \
  b_{f}
  &&& \text{forget gate} \\
\label{eq:mlstm_recurrent_end}
\mathbf{o}_t \ &= \ \sigma \left( \tilde{\mathbf{o}}_t \right) \ , \qquad \qquad \quad \ \ \ \,
  \tilde{\mathbf{o}}_t  \ = \ \mathbf{W}_{\mathbf{o}} \ \mathbf{x}_t \ + \
  \mathbf{b}_{\mathbf{o}}
  &&&\text{output gate} 
\end{align}

The mLSTM architecture supports several memory cells, similar to the standard LSTM. With mLSTM, the use of multiple heads is functionally the same as multiple cells, due to the absence of memory mixing between components. To ensure the stability of mLSTM’s exponential gates, we apply the same stabilization approaches utilized for sLSTM, as indicated by Equation~\ref{eq:slstmstabil}. Because mLSTM does not incorporate memory mixing, its recurrence relation can be restated in a parallel format. Additional discussion can be found in Appendix~\ref{sec:appmLSTM}.

\subsection{xLSTM Architecture}
\label{sec:xLSTMarch}

\paragraph{xLSTM Blocks.}
\label{sec:xLSTMblock}
An xLSTM block aims to capture and encode past information in a high-dimensional non-linear space, facilitating the disambiguation of different past trajectories or contextual histories. Effective differentiation of these histories is critical for predicting subsequent elements in a sequence, such as the upcoming token. We invoke Cover’s Theorem~\citep{Cover:65}, which implies that patterns that are non-linearly mapped to higher-dimensional spaces are more amenable to linear separation compared to those in the original space. In our approach, we examine two types of residual block designs: (i) A residual block employing a post up-projection approach (akin to Transformers), where the summary of historical data is first generated non-linearly within the original representational space, followed by a linear transformation to a higher-dimensional space, application of a non-linear activation, and finally a linear projection returning to the original dimensionality; refer to the left panel of Figure~\ref{fig:Backbones} and the third column in Figure~\ref{fig:xlstm_sketch}. A comprehensive schematic is available in Figure~\ref{fig:appsLSTM_detailed} in the appendix. (ii) A residual block featuring pre up-projection (similar to State Space Models), which first projects input data linearly into a high-dimensional space,

The past is non-linearly encoded within a high-dimensional space before being projected linearly back to the original space. In xLSTM blocks that utilize sLSTM, the post up-projection block is employed. Conversely, for xLSTM blocks with mLSTM, the pre up-projection block is used due to the increased memory capacity afforded by the high-dimensional space. For additional clarification, see the left panel of Figure~\ref{fig:Backbones}, the third column in Figure~\ref{fig:xlstm_sketch}, or consult Figure~\ref{fig:appsLSTM_detailed} provided in the appendix.

\paragraph{xLSTM Architecture. }
\label{sec:xLSTMarchitecure}
The xLSTM architecture is developed by stacking modular components in a residual fashion~\citep{Srivastava:15,He:16}. Our design utilizes the standard pre-LayerNorm~\citep{Ba:16b} residual framework that is prevalent among modern Large Language Models. Refer to the rightmost column of Figure~\ref{fig:xlstm_sketch} for an illustration.

\subsection{Memory and Speed Considerations}

In contrast to Transformers, xLSTM networks feature linear computational complexity and maintain constant memory requirements regardless of sequence length. Due to the compressive nature of xLSTM memory, these networks are particularly advantageous for industrial uses and edge device implementations.

The mLSTM’s memory mechanism is parameter-free; however, it incurs significant computational overhead due to both the $d \times d$ memory matrix and the $d \times d$ update process. In this architecture, we prioritize increased memory capacity at the expense of greater computational demands. Despite this, these operations can be efficiently parallelized on GPUs, resulting in only negligible impact on overall runtime.

Although mLSTM allows for parallel processing similar to approaches like FlashAttention~\citep{Dao:22,Dao:23} or GLA~\citep{Yang:23arxiv}, sLSTM lacks this parallelizability because of its reliance on memory mixing through hidden-to-hidden connections. Nevertheless, we introduced an efficient CUDA implementation, optimizing GPU memory usage down to the register level, resulting in performance that is generally no more than twice as slow as mLSTM.

\section{Related Work}
\label{sec:related}

\paragraph{Linear Attention.} A variety of strategies have been proposed to address the quadratic time complexity of the Transformer with respect to context length, with the aim of achieving linear scaling for attention. The Synthesizer, for instance, generates synthetic attention weights independently of token-to-token dependencies~\citep{Tay:20arxiv}. Linformer introduces a low-rank projection to approximate self-attention in a manner that is linear with respect to sequence length~\citep{Wang:20arxiv}. The Linear Transformer modifies the attention computation so that it operates in linear time~\citep{Katharopoulos:20}. Performer applies a positive orthogonal random feature framework to approximate the softmax operation in attention linearly~\citep{Choromanski:21}. In place of traditional attention, Structured Global Convolution (SGConv)~\citep{Li:22} and the Hyena Hierarchy~\citep{Poli:23} adopt efficient long convolutions for handling long-range dependencies.

\paragraph{State Space Models.} In recent years, State Space Models (SSMs) have gained significant attention due to their linear scalability with respect to sequence length and their competitive results when benchmarked against Transformer models. Among the pioneering efforts in this area is the Structured State Space sequence model (S4)~\citep{Gu:21}. Subsequent advancements include the introduction of the Diagonal State Space (DSS) model~\citep{Gupta:22}, Gated State Space (GSS) models~\citep{Mehta:22}, the S5 model~\citep{Smith:22}, Bidirectional Gated SSM (BiGS)~\citep{Wang:22}, H3 model~\citep{Fu:23}, and Mamba~\citep{Gu:24arxiv}.

\paragraph{Recurrent Neural Networks.} Recurrent Neural Networks (RNNs) have been explored as potential alternatives to Transformer architectures and attention mechanisms because of their linear scaling with respect to sequence length. Language modeling experiments utilizing Deep Linear Recurrent Units (LRUs) have achieved strong performance~\citep{Orvieto:23, De:24arxiv}, and comparable success has been observed with Hierarchically Gated Linear RNN (HGRN) models~\citep{Qin:23} and their successor HGRN2~\citep{Qin:24arxiv}. Among notable large language model architectures based on RNNs is RWKV~\citep{Peng:23arxivshort,Peng:24arxivshort}, which has demonstrated results on par with Transformer-based systems.

\paragraph{Gating.}
The concept of gating, a foundational aspect of LSTM, has been revisited and variously applied in numerous recent models. This mechanism is featured prominently in models such as HGRN~\citep{Qin:23}, HGRN2~\citep{Qin:24arxiv}, Gated Linear Attention (GLA)~\citep{Yang:23arxiv}, Gated State Space (GSS) models~\citep{Mehta:22}, Bidirectional Gated SSM (BiGS)~\citep{Wang:22}, Moving Average Equipped Gated Attention (MEGA)~\citep{Ma:22}, RWKV~\citep{Peng:23arxivshort}, as well as Mamba~\citep{Gu:24arxiv}.

\paragraph{Covariance Update Rule.} In order to improve storage capacity, we augmented the mLSTM cell by incorporating a matrix memory governed by a covariance update rule. Several other models employ similar update mechanisms, including Fast Weight Programmers \citep{Schmidhuber:92ncfastweights,Schlag:21}, RWKV-5 and RWKV-6~\citep{Peng:24arxivshort}, Retention~\citep{Sun:23arxiv}, Linear Transformer~\citep{Katharopoulos:20}, and HGRN2~\citep{Qin:24arxiv}.

\paragraph{Most Related.} The models that are most similar in spirit to xLSTM include Retention~\citep{Sun:23arxiv}, RWKV~\citep{Peng:23arxivshort,Peng:24arxivshort}, and HGRN2~\citep{Qin:24arxiv}. All of these models incorporate elements such as matrix memory and/or gating mechanisms. Nevertheless, unlike the recently introduced sLSTM, these methods do not support memory mixing. The ability to mix memories is vital for addressing state tracking tasks, making LSTMs intrinsically more powerful in this regard than State Space Models (SSMs) and Transformers~\citep{Merrill:24,Deletang:23}. Accurate state tracking, in turn, is essential for tasks such as code interpretation or maintaining entity information throughout lengthy textual sequences.

\paragraph{Residually Stacking Architectures.} Consistent with the design of most modern large-scale deep learning systems, xLSTM models employ residual stacking of modular units~\citep{Srivastava:15,He:16}.

This approach underpins the success of deep convolutional architectures~\citep{He:16} as well as the Transformer framework~\citep{Vaswani:17}. Notably, the Transformer architecture is the foundational technology powering Large Language Models (LLMs) such as GPT-3~\citep{Brown:20short}, ChatGPT~\citep{Schulman:22short}, GPT-4~\citep{Achiam:23gpt4short}, Megatron-LM~\citep{Shoeybi:19}, Gopher~\citep{Rae:21short}, ERNIE 3.0 Titan~\citep{Wang:21arxivshort}, GLaM~\citep{Du:21glamshort}, the Chinese M6~\citep{Lin:21}, multilingual AlexaTM 20B~\citep{Soltan:22}, OPT~\citep{Zhang:22opt}, Chinchilla~\citep{Hoffmann:22short}, BLOOM~\citep{Scao:22arxivshort}, GLM-130B~\citep{Zeng:22arxivshort}, LaMDA~\citep{Thoppilan:22short}, PaLM~\citep{Chowdhery:22short}, Llama~\citep{Touvron:23arxiv}, and Gemini~\citep{GeminiTeam:23arxiv,Reid:24arxiv}.

\section{Experiments}
\label{sec:Exp}

This section presents an empirical analysis of xLSTM in the context of language modeling, highlighting its performance relative to existing approaches. Section~\ref{sec:ExpSyntethic} is dedicated to probing xLSTM's abilities on a suite of synthetic benchmarks. Subsequently, in Section~\ref{sec:ExpComparison}, we benchmark the validation perplexity across various language modeling techniques, all trained with 15B tokens sourced from SlimPajama~\citep{Soboleva:23}. We also conduct ablation experiments on xLSTM using this dataset. Further, we analyze and compare the scaling laws governing the methods, following the methodologies described in \citet{Kaplan:20} and \citet{Brown:20short}. 

Moving to Section~\ref{sec:ExpLanguage}, we deepen the comparative study by conducting comprehensive language modeling experiments. Here, both xLSTM and top-performing alternatives from Section~\ref{sec:ExpComparison} are trained with a substantially larger 300B-token subset of SlimPajama~\citep{Soboleva:23}. Our evaluation unfolds in four stages: assessing the models' capacity to extrapolate to extended context lengths, comparing their validation perplexity and downstream performance as outlined in~\citep{Sutawika:24}, evaluating on all 571 domains included in the PALOMA language benchmark dataset~\citep{Magnusson:23arxivshort}, and finally, re-examining scaling trends when models are exposed to data volumes twenty times greater than in prior experiments.

\label{sec:xlstm_blocks_notation}
In our experiments, we denote xLSTM[$a$:$b$] to signify the proportion $a/b$ of xLSTM blocks that are constructed using mLSTM as opposed to sLSTM modules. To illustrate, xLSTM[7:1] refers to an arrangement where, from a total of eight blocks, seven utilize mLSTM, while one employs sLSTM. When working with a standard configuration of 48 total blocks, this arrangement yields 42 mLSTM-based and 6 sLSTM-based blocks. Additionally, across all experiments, we apply up-projection blocks prior to the mLSTM blocks (“pre”) and after the sLSTM blocks (“post”).

\subsection{Synthetic Tasks and Long Range Arena}
\label{sec:ExpSyntethic}
To begin, we evaluate how well xLSTM's innovative exponential gating combined with memory mixing performs on formal languages~\citep{Deletang:23}. Next, we investigate the capabilities of xLSTM’s novel matrix memory through the Multi-Query Associative Recall task~\citep{Arora:23arxiv}. Lastly, we measure xLSTM's ability to handle extended input sequences within the Long Range Arena benchmark~\citep{Tay:21}.

\paragraph{Evaluation of xLSTM's Exponential Gating with Memory Mixing.} We investigate the capabilities of xLSTM's exponential gating mechanism, which incorporates memory mixing, hypothesized to be crucial for tackling state tracking challenges~\citep{Merrill:24, Merrill:23}. Our approach involves adapting and broadening the suite of formal language tasks introduced by \citet{Deletang:23}, allowing models to be trained on varying sequence lengths in order to assess their generalization to longer lengths. Comprehensive explanations of these tasks, along with expanded empirical results, are provided in Appendix~\ref{sec:appExpSynthetic-formalLang}. 

xLSTM's performance is benchmarked alongside a variety of architectures, including Transformers, State Space Models, and classic Recurrent Neural Networks. For evaluation, we measure accuracy solely on those tokens that are relevant for each respective task, with reported scores normalized between 0 (chance-level performance) and 1 (perfect performance). Specifically, we contrast 2-block variants of each model—namely, xLSTM[0:1] (sLSTM only), xLSTM[1:0] (mLSTM only), xLSTM[1:1], Llama, Mamba, RWKV, Retention, Hyena, LSTM, and LSTM incorporated within Transformer blocks (LSTM (Block)). The experimental outcomes are depicted in Figure~\ref{fig:formal-main}.

Notably, models like standard Transformers or State Space Models that do not employ memory mixing—and thus lack intrinsic state tracking capabilities—fail to handle tasks grounded in regular grammars, such as the parity task. This observation aligns with prior studies indicating that Transformers and State Space Models are inherently less expressive than RNN-based architectures in such settings~\citep{Merrill:24, Merrill:23, Deletang:23}.

\paragraph{Test of xLSTM's Memory Capacities on Associative Recall Tasks.} This study evaluates the matrix-based memory mechanism introduced in xLSTM by measuring its memory capacity using the Multi-Query Associative Recall task~\citep{Arora:23arxiv}. In each test instance, key-value pairs are sampled at random from a comprehensive vocabulary, and models are required to store this information for subsequent recall. To raise the challenge posed by the original task, the number of key-value pairs is increased to as many as 256, and the sequence length is extended up to 2048 tokens. This design allows a more thorough assessment of various architectures' memory abilities. The comparison includes 2-block versions of Llama, Mamba, RWKV-5, RWKV-6, xLSTM[1:1], and xLSTM[1:0]. Model performance is measured by recall accuracy. Transformers such as Llama, known for their exponential memory relative to the representational dimension~\citep{Ramsauer:21}, provide a benchmark for this task. Figure~\ref{fig:mqar-main} reports the findings.
Among non-Transformer architectures, xLSTM[1:1] exhibits the highest performance, even at smaller model scales. Notably, incorporating the sLSTM block does not impede memory capabilities; instead, it appears to enhance them, particularly in the most demanding settings involving 256 key-value pairs. Supplementary analyses can be found in Appendix~\ref{sec:appExpSynthetic-mqar}, including extrapolation studies demonstrating that the advanced memory structure of xLSTM generalizes effectively to longer contexts than those used during its training phase.

\paragraph{Test of xLSTM's Long Context Capabilities on Long Range Arena.} In order to evaluate xLSTM’s ability to process lengthy sequences and extensive contexts, we conducted comparisons with alternative approaches on the Long Range Arena~\citep{Tay:21}. Across every task, xLSTM delivers robust and reliable results, indicating that its architectural design is highly effective for managing diverse challenges associated with long context modeling.

For more details, see Appendix~\ref{sec:appExpSynthetic-lra}.

\subsection{Method Comparison and Ablation Study}
\label{sec:ExpComparison}

In order to investigate the core research question—namely, the potential of our novel LSTM variants when applied to large-scale language modelling—we conduct experiments involving xLSTMs, Transformers, State Space Models, and additional approaches, all trained under consistent auto-regressive conditions using 15B tokens from SlimPajama. We evaluate model performance on a validation dataset and conduct ablation analyses specifically for xLSTMs.

\paragraph{Comparative Evaluation of xLSTM Against Competing Models.} In our evaluation, models are trained using a dataset comprising 15B tokens sourced from SlimPajama~\citep{Soboleva:23}. Performance is measured by computing perplexity on the validation subset. Our analysis includes several approaches: xLSTM (the proposed architecture), GPT-3 (Transformer-based)~\citep{Brown:20short}, Llama (Transformer-based)~\citep{Touvron:23arxiv}, H3 (state space model, SSM)~\citep{Fu:23}, Mamba (SSM)~\citep{Gu:24arxiv}, RWKV-4/5/6 (RNN-based)~\citep{Peng:23arxivshort,Peng:24arxivshort}, GLA (a linear Transformer)~\citep{Yang:23arxiv}, HGRN and HGRN2 (RNN-based)~\citep{Qin:23,Qin:24arxiv}, as well as RetNet (linear Transformer)~\citep{Sun:23arxiv}, Hyena (linear Transformer)~\citep{Poli:23}, and xLSTM variants xLSTM[1:0] and xLSTM[7:1] (notations described in Section~\ref{sec:xlstm_blocks_notation}). 
The training utilized mixed precision; however, for RWKV-5, RWKV-6, GLA, and HGRN2, automated mixed precision from PyTorch was not employed (refer to Appendix Section~\ref{sec:appGenTrainProc} for further clarifications).
For organization, the methods are grouped into three categories: (a) Transformer architectures, (b) State Space Models (SSMs), and (c) Recurrent Neural Networks (RNNs) alongside linear Transformer variants. Linear Transformers implement a linearized alternative to standard Transformer attention. All model configurations are scaled to correspond to a GPT-3 baseline with 350M parameters, that is, an embedding dimension of 1024 and 24 residual layers. GPT-3 exclusively employs shared token and output embeddings, leading to a reduction in parameter count relative to others. 
According to Table~\ref{tab:spaj15b_model_results}, xLSTM achieves the lowest validation perplexity, surpassing all compared baselines. Extended results are provided in Appendix~\ref{sec:appExpComparison}. 
Furthermore, scaling dynamics for the experimental setup, depicted in Figure~\ref{fig:temp_sclaw15B}, suggest that xLSTM is expected to maintain its superior performance as model size increases.

\paragraph{Ablation Studies.}
As shown in Table~\ref{tab:spaj15b_model_results} and Figure~\ref{fig:temp_sclaw15B}, xLSTM attains strong performance in language modeling tasks after training on 15B SlimPajama tokens. To investigate the specific contributions of each architectural modification from LSTM to xLSTM, we incrementally transform a standard LSTM model toward xLSTM. The process begins by incorporating LSTM layers into pre-LayerNorm residual backbones. The second modification introduces a post up-projection block. In the final step, we supplement the architecture with exponential gating and matrix memory components.

The results are shown in Table~\ref{tab:ablstudies} (top). 

Ablation experiments suggest that the observed performance gains result from the combined effect of the exponential gating mechanism and the matrix memory. Given the central role of gating within RNNs and State Space Models, we also assess a range of gating variants. As indicated in Table~\ref{tab:ablstudies} (bottom), we find that when each gate is both learnable and modulated by the input, it contributes progressively to performance improvements. Further analyses regarding each individual backbone component can be found in Appendix~\ref{sec:appAblDetails}.

\subsection{xLSTM as Large Language Model}
\label{sec:ExpLanguage}

To conclude our investigation, we conduct extensive language modeling experiments aimed at evaluating xLSTM’s viability as a large language model (LLM). To this end, we scale up the dataset size, utilizing 300B tokens from SlimPajama for training. Notably, this token count aligns with the setups employed in models such as Mamba~\citep{Gu:24arxiv} and Griffin~\citep{De:24arxiv}.

We benchmark xLSTM against RWKV-4, Llama, and Mamba, each representing a distinct methodological category as outlined in Section~\ref{sec:ExpComparison}. RWKV-4 is chosen to represent the RNN class, as suitable training precision parameters for RWKV-5, RWKV-6, and HGRN2 were established only after the commencement of the 300B token experiments (refer to Appendix~\ref{sec:appGenTrainProc} for details). Our experiments involve a range of model sizes—125M, 350M, 760M, and 1.3B. All models are evaluated for their ability to extrapolate to longer sequence lengths and assessed on the validation set for general performance. In addition, we analyze their effectiveness on a variety of downstream tasks, measure their language modeling performance across 571 text domains from the PALOMA benchmark, and study their scaling law properties.

\paragraph{Sequence Length Extrapolation.}
\label{sec:spaj300B_ctx_extrapolation}
We begin by evaluating the ability of large models—specifically, those with 1.3B parameters—namely xLSTM, RWKV-4, Llama, and Mamba, to extrapolate to longer sequence lengths. Although all these models are trained using a context length of 2048, we examine their performance at extended context lengths, reaching up to 16384. The corresponding outcomes are displayed in Figure~\ref{fig:spaj300B_seq_extrapolation}. Notably, among the various approaches, xLSTM models consistently achieve low perplexity values across these longer contexts.

\paragraph{Validation Perplexity and Downstream Tasks.}
\label{sec:spaj_downstream_eval}
Additionally, across every model scale, we assess xLSTM, RWKV-4, Llama, and Mamba models by measuring their next token prediction accuracy on the SlimPajama validation dataset, as well as by evaluating their effectiveness on downstream tasks aimed at testing common sense reasoning capabilities.

The third column in Table~\ref{tab:lmeval} presents the validation set perplexity achieved by each method. Notably, xLSTM[1:0] and xLSTM[7:1] outperform other models in terms of validation set perplexity across all tested model sizes. The remaining columns of Table~\ref{tab:lmeval} show the results on downstream tasks. Across nearly all tasks and model sizes, xLSTM consistently yields the best performance—exceptions arise only on the ARC task, where Mamba is occasionally superior. Further specifics can be found in Appendix~\ref{sec:appExpLanguage}.

\paragraph{Performance on PALOMA Language Tasks.} Third, we evaluate the capabilities of xLSTM, RWKV-4, Llama, and Mamba in next token prediction across all model sizes on PALOMA language tasks~\citep{Magnusson:23arxivshort}. To assess model accuracy, we report perplexity scores for predicting the subsequent token across 571 varied text domains, encompassing sources such as nytimes.com and Reddit's r/depression. Table~\ref{tab:ppleval} presents these perplexity results, with columns separating standard language modeling domains (the first seven) from the more detailed, domain-specific tasks (the final five). On these tasks, xLSTM[1:0] consistently surpasses xLSTM[7:1] in performance.

xLSTM[1:0] achieves lower perplexity compared to Mamba in 568 of 571 text domains (99.5\%), outperforms Llama in 486 out of 571 cases (85.1\%), and demonstrates lower perplexity than RWKV-4 in 570 of 571 domains (99.8\%). Further information is provided in Appendix~\ref{sec:appExpLanguage}.

\paragraph{Scaling Laws.} Fourth, we investigate the presence of power-law scaling, which enables projection of model performance at increased parameter counts~\citep{Kaplan:20,Brown:20short}. The scaling results are illustrated in Figure~\ref{fig:temp_sclaw300B}. While all evaluated models follow comparable scaling trends, their absolute performance levels are offset. Specifically, RWKV-4 demonstrates the lowest performance, with Llama and Mamba following in succession. Notably, xLSTM outperforms Mamba by a margin that mirrors the performance gap between Mamba and Llama. This scaling pattern suggests that as models scale up, xLSTM is expected to maintain its comparative advantage over Transformer-based and State-Space architectures.

\textbf{Generation Times and Maximal Throughput.} In our final set of experiments, we assess the text generation latency as depicted in Figure~\ref{fig:inference_time_generation} and evaluate peak throughput in Figure~\ref{fig:inference_time_decoding_throughput} (left) for various xLSTM configurations at the 1.3B parameter scale. We include comparisons to HuggingFace implementations of Mamba, Llama, and RWKV models of comparable size, employing a static key-value cache for the Llama variant. During these experiments, we observed that full cache compilation for the Transformer architecture and compiling the Mamba model using \texttt{torch.compile} were not functional. All generative evaluations utilized a batch size of 1 with a pre-fill value of 16, a setting most advantageous to the Transformer approach.

Figure~\ref{fig:inference_time_generation} illustrates that xLSTM, alongside other recurrent models such as Mamba and RWKV-4, exhibits linear scaling in generation time, whereas Llama demonstrates quadratic scaling behavior. For throughput analysis, we tested the Llama model under a variety of batch size and pre-fill conditions. Notably, Figure~\ref{fig:inference_time_decoding_throughput} (right) demonstrates that xLSTM is capable of supporting significantly larger batch sizes relative to Llama, owing to xLSTM’s constant memory requirements, resulting in superior throughput performance.

\section{Limitations}

(i) Unlike mLSTM, sLSTM’s approach to memory mixing inherently restricts parallel computation, rendering fast parallel implementations infeasible. However, we have engineered an efficient CUDA kernel for sLSTM that is currently less than twice as slow as our parallelized mLSTM version. (ii) Our mLSTM CUDA kernels are not yet fully optimized, resulting in execution speeds roughly four times slower compared to FlashAttention or the scanning procedure applied in Mamba. Employing techniques similar to those in FlashAttention may yield substantial speedups with further kernel optimization. (iii) The mLSTM’s matrix memory necessitates processing $d \times d$ matrices, leading to increased computational demands. Nevertheless, because the processes of memory updating and retrieval are parameter-free and parallelizable via standard matrix algebra, the actual increase in computation time from the intricate memory structure remains minimal. (iv) Careful selection of initial values for the forget gates is required. (v) Since the size of matrix memory does not vary with sequence length, expanding sequence length can risk memory overload at larger context sizes; however, this has not presented practical issues for contexts as large as 16k tokens, as discussed in Section~\ref{sec:spaj300B_ctx_extrapolation}. (vi) Owing to the considerable computational requirements of large-scale language model experimentation, we did not conduct comprehensive optimization of either architectural choices or hyperparameters—especially for more sizeable xLSTM variants. We expect that rigorous and targeted optimization will be necessary for xLSTM models to achieve maximal performance.

\section{Conclusion}
Our investigation provides a partial response to our initial inquiry: To what extent does the performance of language modeling improve by scaling LSTM architectures to the order of billions of parameters? Presently, we can state: ``At least to the level achieved by leading architectures such as Transformers and State Space Models.'' We have proposed an augmented variant, xLSTM, which incorporates exponential gating mechanisms, memory mixing, and a revised memory architecture. When benchmarked against advanced frameworks, including Transformers and State Space Models, xLSTM demonstrates competitive results on language modeling benchmarks. Observations from the scaling laws suggest that expanding xLSTM to greater model sizes positions it as a viable contender alongside today's dominant Large Language Models that are Transformer-based. Moreover, xLSTM shows promise for substantial advancements in diverse areas such as Reinforcement Learning, Time Series Forecasting, and physical system modeling.

\end{document}